\section{Eigenvalue Problems for \texorpdfstring{$y''+\lambda y=0$}{y''+lambda y=0}}

\noindent
In Chapter~12 we'll study partial differential equations that arise in problems of heat conduction, wave propagation, and potential theory. The purpose of this chapter is to develop tools required to solve these equations. In this section we consider the following problems, where $\lambda$ is a real number and $L>0$:

Problem 1: \hskip5em $y''+\lambda y=0,\quad y(0)=0,\quad y(L)=0$
\medskip

Problem 2: \hskip5em $y''+\lambda y=0,\quad y'(0)=0,\quad y'(L)=0$
\medskip

Problem 3: \hskip5em $y''+\lambda y=0,\quad y(0)=0,\quad y'(L)=0$
\medskip

Problem 4: \hskip5em $y''+\lambda y=0,\quad y'(0)=0,\quad y(L)=0$
\medskip

Problem 5: \hskip5em $y''+\lambda y=0,\quad y(-L)=y(L), \quad y'(-L)=y'(L)$
\medskip

In each problem the conditions following the differential equation are called {\color{blue}\it boundary conditions\/}. Note that the boundary conditions in Problem~5, unlike those in Problems~1-4, don't require that $y$ or $y'$ be zero at the boundary points, but only that $y$ have the same value at $x=\pm L$ , and that $y'$ have the same value at $x=\pm L$. We say that the boundary conditions in Problem~5 are {\color{blue}\it periodic\/}.

Obviously, $y\equiv0$ (the trivial solution) is a solution of Problems~1-5 for any value of $\lambda$. For most values of $\lambda$, there are no other solutions. The interesting question is this:

{\color{blue}\it For what values of $\lambda$ does the problem have nontrivial solutions, and what are they?}

A value of $\lambda$ for which the problem has a nontrivial solution is an {\color{blue}\it eigenvalue\/} of the problem, and the nontrivial solutions are $\lambda$-{\color{blue}\it eigenfunctions\/}, or {\color{blue}\it eigenfunctions associated with $\lambda$\/}. Note that a nonzero constant multiple of a $\lambda$-eigenfunction is again a $\lambda$-eigenfunction.

Problems~1-5 are called {\color{blue}\it eigenvalue problems\/}. {\color{blue}\it Solving\/} an eigenvalue problem means finding all its eigenvalues and associated eigenfunctions. We'll take it as given here that all the eigenvalues of Problems~1-5 are real numbers. This is proved in a more general setting in Section~13.2.

\begin{theorem}
  \color{blue}
  \label{thmtype:11.1.1} Problems $1$--$5$ have no negative eigenvalues. Moreover, $\lambda=0$ is an eigenvalue of Problems $2$ and $5$, with associated eigenfunction $y_{0}=1$, but $\lambda=0$ isn't an eigenvalue of Problems~$1$, $3$, or $4$.
\end{theorem}

\proof We consider Problems 1-4, and leave Problem 5 to you (Exercise~\ref{exer:11.1.1}). If $y''+\lambda y=0$, then $y(y''+\lambda y)=0$, so
\[
  \int_{0}^{L}y(x)(y''(x)+\lambda y(x))\,dx=0;
\]
therefore,
\begin{equation}
  \label{eq:11.1.1}\lambda\int_{0}^{L}y^{2}(x)\,dx=-\int_{0}^{L}y(x)y''(x)\,dx.
\end{equation}
Integration by parts yields
\begin{equation}
  \label{eq:11.1.2}
  \begin{array}{rcl}
    \dst\int_{0}^{L}y(x)y''(x)\,dx & = & \dst y(x)y'(x)\lims0L-\int_{0}^{L} (y'(x))^{2}\,dx    \\
                                   & = & \dst y(L)y'(L)-y(0)y'(0)-\int_{0}^{L}(y'(x))^{2}\,dx.
  \end{array}
\end{equation}
However, if $y$ satisfies any of the boundary conditions of Problems 1-4, then
\[
  y(L)y'(L)-y(0)y'(0)=0;
\]
hence, \eqref{eq:11.1.1} and \eqref{eq:11.1.2} imply that
\[
  \lambda\int_{0}^{L}y^{2}(x)\,dx=\int_{0}^{L}(y'(x))^{2}\,dx.
\]
If $y\not\equiv0$, then $\int_{0}^{L}y^{2}(x)\,dx>0$. Therefore $\lambda\ge0$ and, if $\lambda=0$, then $y'(x)=0$ for all $x$ in $(0,L)$ (why?), and $y$ is constant on $(0,L)$. Any constant function satisfies the boundary conditions of Problem~2, so $\lambda=0$ is an eigenvalue of Problem~2 and any nonzero constant function is an associated eigenfunction. However, the only constant function that satisfies the boundary conditions of Problems $1$, $3$, or $4$ is $y\equiv0$. Therefore $\lambda=0$ isn't an eigenvalue of any of these problems.

\begin{example}[Problem 1]
  \label{example:11.1.1} \rm Solve the eigenvalue problem
  \begin{equation}
    \label{eq:11.1.3}y''+\lambda y=0,\quad y(0)=0,\quad y(L)=0.
  \end{equation}

  \solution From Theorem~\ref{thmtype:11.1.1}, any eigenvalues of \eqref{eq:11.1.3} must be positive. If $y$ satisfies \eqref{eq:11.1.3} with $\lambda>0$, then
  \[
    y=c_{1}\cos\sqrt{\lambda}\,x+c_{2}\sin\sqrt{\lambda}\,x,
  \]
  where $c_{1}$ and $c_{2}$ are constants. The boundary condition $y(0)=0$ implies that $c_{1}=0$. Therefore $y=c_{2}\sin\sqrt{\lambda}\,x$. Now the boundary condition $y(L)=0$ implies that $c_{2}\sin\sqrt{\lambda}\,L=0$. To make $c_{2}\sin\sqrt{\lambda}\,L=0$ with $c_{2}\ne0$, we must choose $\sqrt{\lambda}=n\pi/L$, where $n$ is a positive integer. Therefore $\lambda_{n}=n^{2}\pi^{2}/L^{2}$ is an eigenvalue and
  \[
    y_{n}=\sin\frac{n\pi x}{L}
  \]
  is an associated eigenfunction. \bbox
\end{example}

For future reference, we state the result of Example~\ref{example:11.1.1} as a theorem.

\begin{theorem}
  \color{blue}
  \label{thmtype:11.1.2} The eigenvalue problem
  \[
    y''+\lambda y=0,\quad y(0)=0,\quad y(L)=0
  \]
  has infinitely many positive eigenvalues $\lambda_{n}=n^{2}\pi^{2}/L^{2}$, with associated eigenfunctions
  \[
    y_{n}=\sin\frac{n\pi x}{L},\quad n=1,2,3,\dots.
  \]
  There are no other eigenvalues.
\end{theorem}

We leave it to you to prove the next theorem about Problem~2 by an argument like that of Example~\ref{example:11.1.1} (Exercise~\ref{exer:11.1.17}).

\begin{theorem}
  \color{blue}
  \label{thmtype:11.1.3} The eigenvalue problem
  \[
    y''+\lambda y=0,\quad y'(0)=0,\quad y'(L)=0
  \]
  has the eigenvalue $\lambda_{0}=0$, with associated eigenfunction $y_{0}=1$, and infinitely many positive eigenvalues $\lambda_{n}=n^{2}\pi^{2}/L^{2}$, with associated eigenfunctions
  \[
    y_{n}=\cos\frac{n\pi x}{L},n=1,2,3\dots.
  \]
  There are no other eigenvalues.
\end{theorem}

\begin{example}[Problem 3]
  \label{example:11.1.2} \rm Solve the eigenvalue problem
  \begin{equation}
    \label{eq:11.1.4}y''+\lambda y=0,\quad y(0)=0,\quad y'(L)=0.
  \end{equation}

  \solution From Theorem~\ref{thmtype:11.1.1}, any eigenvalues of \eqref{eq:11.1.4} must be positive. If $y$ satisfies \eqref{eq:11.1.4} with $\lambda>0$, then
  \[
    y=c_{1}\cos\sqrt{\lambda}\,x+c_{2}\sin\sqrt{\lambda}\,x,
  \]
  where $c_{1}$ and $c_{2}$ are constants. The boundary condition $y(0)=0$ implies that $c_{1}=0$. Therefore $y=c_{2}\sin\sqrt{\lambda}\,x$. Hence, $y'=c_{2}\sqrt{\lambda}\cos\sqrt{\lambda}\,x$ and the boundary condition $y'(L)=0$ implies that $c_{2}\cos\sqrt{\lambda}\,L=0$. To make $c_{2}\cos\sqrt{\lambda}\,L=0$ with $c_{2}\ne0$ we must choose
  \[
    \sqrt{\lambda}=\frac{(2n-1)\pi}{2L},
  \]
  where $n$ is a positive integer. Then $\lambda_{n}=(2n-1)^{2}\pi^{2}/4L^{2}$ is an eigenvalue and
  \[
    y_{n}=\sin\frac{(2n-1)\pi x}{2L}
  \]
  is an associated eigenfunction. \bbox
\end{example}

For future reference, we state the result of Example~\ref{example:11.1.2} as a theorem.

\begin{theorem}
  \color{blue}
  \label{thmtype:11.1.4} The eigenvalue problem
  \[
    y''+\lambda y=0,\quad y(0)=0,\quad y'(L)=0
  \]
  has infinitely many positive eigenvalues $\lambda_{n}=(2n-1)^{2}\pi^{2}/4L^{2}$, with associated eigenfunctions
  \[
    y_{n}=\sin\frac{(2n-1)\pi x}{2L},\quad n=1,2,3,\dots.
  \]
  There are no other eigenvalues.
\end{theorem}

We leave it to you to prove the next theorem about Problem~4 by an argument like that of Example~\ref{example:11.1.2} (Exercise~\ref{exer:11.1.18}).

\begin{theorem}
  \color{blue}
  \label{thmtype:11.1.5} The eigenvalue problem
  \[
    y''+\lambda y=0,\quad y'(0)=0,\quad y(L)=0
  \]
  has infinitely many positive eigenvalues $\lambda_{n}=(2n-1)^{2}\pi^{2}/4L^{2}$, with associated eigenfunctions
  \[
    y_{n}=\cos\frac{(2n-1)\pi x}{2L},\quad n=1,2,3,\dots.
  \]
  There are no other eigenvalues.
\end{theorem}

\begin{example}[Problem 5]
  \label{example:11.1.3} \rm Solve the eigenvalue problem
  \begin{equation}
    \label{eq:11.1.5}y''+\lambda y=0,\quad y(-L)=y(L),\quad y'(-L)=y'(L).
  \end{equation}

  \solution From Theorem~\ref{thmtype:11.1.1}, $\lambda=0$ is an eigenvalue of \eqref{eq:11.1.5} with associated eigenfunction $y_{0}=1$, and any other eigenvalues must be positive. If $y$ satisfies \eqref{eq:11.1.5} with $\lambda>0$, then
  \begin{equation}
    \label{eq:11.1.6}y=c_{1}\cos\sqrt{\lambda}\,x+c_{2}\sin\sqrt{\lambda}\,x,
  \end{equation}
  where $c_{1}$ and $c_{2}$ are constants. The boundary condition $y(-L)=y(L)$ implies that
  \begin{equation}
    \label{eq:11.1.7}c_{1}\cos(-\sqrt{\lambda}\,L)+c_{2}\sin(-\sqrt{\lambda}\,L)=c_{1}\cos\sqrt{\lambda}\,L+c_{2}\sin\sqrt{\lambda}\,L.
  \end{equation}
  Since
  \begin{equation}
    \label{eq:11.1.8}\cos(-\sqrt{\lambda}\,L)=\cos\sqrt{\lambda}\,L\mbox{\quad and \quad}\sin(-\sqrt{\lambda}\,L)=-\sin\sqrt{\lambda}\,L,
  \end{equation}
  \eqref{eq:11.1.7} implies that
  \begin{equation}
    \label{eq:11.1.9}c_{2}\sin\sqrt{\lambda}\,L=0.
  \end{equation}
  Differentiating \eqref{eq:11.1.6} yields
  \[
    y'=\sqrt{\lambda}\left(-c_{1}\sin\sqrt{\lambda}x+c_{2}\cos\sqrt{\lambda}x\right).
  \]
  The boundary condition $y'(-L)=y'(L)$ implies that
  \[
    -c_{1}\sin(-\sqrt{\lambda}\,L)+c_{2}\cos(-\sqrt{\lambda}\,L)=-c_{1}\sin\sqrt{\lambda}\,L+c_{2}\cos\sqrt{\lambda}\,L,
  \]
  and \eqref{eq:11.1.8} implies that
  \begin{equation}
    \label{eq:11.1.10}c_{1}\sin\sqrt{\lambda}\,L=0.
  \end{equation}

  Eqns.~\eqref{eq:11.1.9} and \eqref{eq:11.1.10} imply that $c_{1}=c_{2}=0$ unless $\sqrt{\lambda}=n\pi/L$, where $n$ is a positive integer. In this case \eqref{eq:11.1.9} and \eqref{eq:11.1.10} both hold for arbitrary $c_{1}$ and $c_{2}$. The eigenvalue determined in this way is $\lambda_{n}=n^{2}\pi^{2}/L^{2}$, and each such eigenvalue has the linearly independent associated eigenfunctions
  \[
    \cos\frac{n\pi x}{L}\mbox{\quad and \quad}\sin\frac{n\pi x}{L}. \eqno{\bbox}
  \]
\end{example}

For future reference we state the result of Example~\ref{example:11.1.3} as a theorem.

\begin{theorem}
  \color{blue}
  \label{thmtype:11.1.6} The eigenvalue problem
  \[
    y''+\lambda y=0,\quad y(-L)=y(L),\quad y'(-L)=y'(L),
  \]
  has the eigenvalue $\lambda_{0}=0$, with associated eigenfunction $y_{0}=1$ and infinitely many positive eigenvalues $\lambda_{n}=n^{2}\pi^{2}/L^{2}$, with associated eigenfunctions
  \[
    y_{1n}=\cos\frac{n\pi x}{L}\mbox{\quad and \quad}y_{2n}=\sin\frac{n\pi x}{L},\quad n=1,2,3,\dots.
  \]
  There are no other eigenvalues.
\end{theorem}

\boxit{Orthogonality}

\noindent
We say that two integrable functions $f$ and $g$ are {\color{blue}\it orthogonal\/} on an interval $[a,b]$ if
\[
  \int_{a}^{b}f(x)g(x)\,dx=0.
\]
More generally, we say that the functions $\phi_{1}$, $\phi_{2}$, \dots, $\phi_{n}$, \dots (finitely or infinitely many) are orthogonal on $[a,b]$ if
\[
  \int_{a}^{b}\phi_{i}(x)\phi_{j}(x)\,dx=0\mbox{\quad whenever \quad}i\ne j.
\]
The importance of orthogonality will become clear when we study Fourier series in the next two sections.

\begin{example}
  \label{example:11.1.4} \rm Show that the eigenfunctions
  \begin{equation}
    \label{eq:11.1.11}1,\,\cos\frac{\pi x}{L},\,\sin\frac{\pi x}{L}, \,\cos\frac{2\pi x}{L},\,\sin\frac{2\pi x}{L},\dots, \cos\frac{n\pi x}{L},\,\sin\frac{n\pi x}{L},\dots
  \end{equation}
  of Problem 5 are orthogonal on $[-L,L]$.
\end{example}

\solution We must show that
\begin{equation}
  \label{eq:11.1.12}\int_{-L}^{L}f(x)g(x)\,dx=0
\end{equation}
whenever $f$ and $g$ are distinct functions from \eqref{eq:11.1.11}. If $r$ is any nonzero integer, then
\begin{equation}
  \label{eq:11.1.13}\int_{-L}^{L}\cos\frac{r\pi x}{L}\,dx =\frac{L}{r\pi}\sin\frac{r\pi x}{L}\bigg|_{-L}^{L}=0.
\end{equation}
and
\[
  \int_{-L}^{L}\sin\frac{r\pi x}{L}\,dx =-\frac{L}{r\pi}\cos\frac{r\pi x}{L}\bigg|_{-L}^{L}=0.
\]
Therefore \eqref{eq:11.1.12} holds if $f\equiv1$ and $g$ is any other function in \eqref{eq:11.1.11}.

If $f(x)=\cos m\pi x/L$ and $g(x)=\cos n\pi x/L$ where $m$ and $n$ are distinct positive integers, then
\begin{equation}
  \label{eq:11.1.14}\int_{-L}^{L}f(x)g(x)\,dx=\int_{-L}^{L}\cos\frac{m\pi x}{L}\cos\frac{n\pi x}{L}\,dx.
\end{equation}
To evaluate this integral, we use the identity
\[
  \cos A\cos B=\frac{1}{2}[\cos(A-B)+\cos(A+B)]
\]
with $A=m\pi x/L$ and $B=n\pi x/L$. Then \eqref{eq:11.1.14} becomes
\[
  \int_{-L}^{L}f(x)g(x)\,dx=\frac{1}{2}\left[\int_{-L}^{L}\cos\frac{(m-n)\pi x}{L}\,dx +\int_{-L}^{L}\cos\frac{(m+n)\pi x}{L}\,dx\right].
\]
Since $m-n$ and $m+n$ are both nonzero integers, \eqref{eq:11.1.13} implies that the integrals on the right are both zero. Therefore \eqref{eq:11.1.12} is true in this case.

If $f(x)=\sin m\pi x/L$ and $g(x)=\sin n\pi x/L$ where $m$ and $n$ are distinct positive integers, then
\begin{equation}
  \label{eq:11.1.15}\int_{-L}^{L}f(x)g(x)\,dx=\int_{-L}^{L}\sin\frac{m\pi x}{L}\sin\frac{n\pi x}{L}\,dx.
\end{equation}
To evaluate this integral, we use the identity
\[
  \sin A\sin B=\frac{1}{2}[\cos(A-B)-\cos(A+B)]
\]
with $A=m\pi x/L$ and $B=n\pi x/L$. Then \eqref{eq:11.1.15} becomes
\[
  \int_{-L}^{L}f(x)g(x)\,dx=\frac{1}{2}\left[\int_{-L}^{L}\cos\frac{(m-n)\pi x}{L}\,dx -\int_{-L}^{L}\cos\frac{(m+n)\pi x}{L}\,dx\right]=0.
\]

If $f(x)=\sin m\pi x/L$ and $g(x)=\cos n\pi x/L$ where $m$ and $n$ are positive integers (not necessarily distinct), then
\[
  \int_{-L}^{L}f(x)g(x)\,dx=\int_{-L}^{L}\sin\frac{m\pi x}{L}\cos\frac{n\pi x}{L}\,dx=0
\]
because the integrand is an odd function and the limits are symmetric about $x=0$. \bbox

Exercises~\ref{exer:11.1.19}-\ref{exer:11.1.22} ask you to verify that the eigenfunctions of Problems~1-4 are orthogonal on $[0,L]$. However, this also follows from a general theorem that we'll prove in Chapter~13.

%\end{comment}

\exercises

\begin{exerciselist}
  \item\label{exer:11.1.1} Prove that $\lambda=0$ is an eigenvalue of Problem~5 with associated eigenfunction $y_{0}=1$, and that any other eigenvalues must be positive. \hint{See the proof of Theorem~\ref{thmtype:11.1.1}.}

  \exercisetext{In Exercises 2-16 solve the eigenvalue problem.}

  \item\label{exer:11.1.2} $y''+\lambda y=0,\quad y(0)=0,\quad y(\pi)=0$

  \item\label{exer:11.1.3} $y''+\lambda y=0,\quad y'(0)=0,\quad y'(\pi)=0$

  \item\label{exer:11.1.4} $y''+\lambda y=0$,\quad $y(0)=0$,\quad $y'(\pi)=0$

  \item\label{exer:11.1.5} $y''+\lambda y=0$,\quad $y'(0)=0$,\quad $y(\pi)=0$

  \item\label{exer:11.1.6} $y''+\lambda y=0,\quad y(-\pi)=y(\pi),\quad y'(-\pi)=y'(\pi)$

  \item\label{exer:11.1.7} $y''+\lambda y=0,\quad y'(0)=0,\quad y'(1)=0$

  \item\label{exer:11.1.8} $y''+\lambda y=0$,\quad $y'(0)=0$,\quad $y(1)=0$

  \item\label{exer:11.1.9} $y''+\lambda y=0,\quad y(0)=0,\quad y(1)=0$

  \item\label{exer:11.1.10} $y''+\lambda y=0,\quad y(-1)=y(1),\quad y'(-1)=y'(1)$

  \item\label{exer:11.1.11} $y''+\lambda y=0$,\quad $y(0)=0$,\quad $y'(1)=0$

  \item\label{exer:11.1.12} $y''+\lambda y=0,\quad y(-2)=y(2),\quad y'(-2)=y'(2)$

  \item\label{exer:11.1.13} $y''+\lambda y=0,\quad y(0)=0,\quad y(2)=0$

  \item\label{exer:11.1.14} $y''+\lambda y=0$,\quad $y'(0)=0$,\quad $y(3)=0$

  \item\label{exer:11.1.15} $y''+\lambda y=0$,\quad $y(0)=0$,\quad $y'(1/2)=0$

  \item\label{exer:11.1.16} $y''+\lambda y=0,\quad y'(0)=0,\quad y'(5)=0$

  \item\label{exer:11.1.17} Prove Theorem~\ref{thmtype:11.1.3}.

  \item\label{exer:11.1.18} Prove Theorem~\ref{thmtype:11.1.5}.

  \item\label{exer:11.1.19} Verify that the eigenfunctions
  \[
    \sin\frac{\pi x}{L}, \,\sin\frac{2\pi x}{L},\dots, \,\sin\frac{n\pi x}{L},\dots\,
  \]
  of Problem~1 are orthogonal on $[0,L]$.

  \item\label{exer:11.1.20} Verify that the eigenfunctions
  \[
    1,\,\cos\frac{\pi x}{L}, \,\cos\frac{2\pi x}{L},\dots, \,\cos\frac{n\pi x}{L},\dots\,
  \]
  of Problem~2 are orthogonal on $[0,L]$.

  \item\label{exer:11.1.21} Verify that the eigenfunctions
  \[
    \sin\frac{\pi x}{2L}, \,\sin\frac{3\pi x}{2L},\dots, \,\sin\frac{(2n-1)\pi x}{2L},\dots\,
  \]
  of Problem~3 are orthogonal on $[0,L]$.

  \item\label{exer:11.1.22} Verify that the eigenfunctions
  \[
    \cos\frac{\pi x}{2L}, \,\cos\frac{3\pi x}{2L},\dots, \,\cos\frac{(2n-1)\pi x}{2L},\dots\,
  \]
  of Problem~4 are orthogonal on $[0,L]$.

  \exercisetext{In Exercises~\ref{exer:11.1.23}-\ref{exer:11.1.26} solve the eigenvalue problem.}

  \item\label{exer:11.1.23} $y''+\lambda y=0,\quad y(0)=0,\quad\dst\int_{0}^{L}y(x)\,dx=0$

  \item\label{exer:11.1.24} $y''+\lambda y=0,\quad y'(0)=0,\quad\dst\int_{0}^{L}y(x)\,dx=0$

  \item\label{exer:11.1.25} $y''+\lambda y=0,\quad y(L)=0,\quad\dst\int_{0}^{L}y(x)\,dx=0$

  \item\label{exer:11.1.26} $y''+\lambda y=0,\quad y'(L)=0,\quad\dst\int_{0}^{L}y(x)\,dx=0$
\end{exerciselist}
